% TOP_PROBLEM: {"id": 30006851, "title": "Marked Length Spectrum Rigidity Conjecture", "category_id": 6, "category": {"id": 6, "name": "geometry", "display_name": "Geometry", "description": "Euclidean and non-Euclidean geometry, geometric structures.", "slug": "geometry", "order_index": 6, "created_at": "2026-07-31T15:26:25.671Z"}, "status": "open"}
% FIRST_POSTED: 2026-09-27
% !TeX program = lualatex
% ATTEMPT_STATUS: solved
% Model: GPT-6 Astra Ultra; continuation dated 27 September 2026.
\documentclass[11pt]{article}
\usepackage[margin=25mm]{geometry}
\usepackage{fontspec}
\setmainfont{Latin Modern Roman}
\usepackage{amsmath,amssymb,mathtools}
\usepackage{unicode-math}
\setmathfont{Latin Modern Math}
\usepackage{xurl,hyperref}
\hypersetup{colorlinks=true,urlcolor=blue}
\setcounter{secnumdepth}{0}
\setlength{\parindent}{0pt}
\setlength{\parskip}{5pt}
\setlength{\emergencystretch}{3em}
\begin{document}
\section*{\texorpdfstring{Marked Length Spectrum Rigidity Conjecture}{Marked Length Spectrum Rigidity Conjecture}}\label{30006851}
\subsection{Definitions and mathematical statement}
Let $M$ be a closed manifold and let $g,h$ be negatively curved Riemannian metrics. Each nontrivial free homotopy class $[\gamma]$ of loops has a unique closed geodesic representative for each metric. Its length defines the \emph{marked length spectrum} $\ell_g([\gamma])$. The conjecture asserts
\[
 \ell_g([\gamma])=\ell_h([\gamma])\text{ for every }[\gamma]
 \quad\Longrightarrow\quad
 \exists f:M\to M\text{ isotopic to }\mathrm{id},\quad f^*h=g.
\]
The marking identifies classes in the same manifold, not just an unordered collection of numerical lengths. The isotopy requirement is part of the supplied target and should not be dropped in favour of an arbitrary isometry. Negative curvature refers to sectional curvature in every tangent two-plane.
\par\smallskip\textit{Formulation and concept references:} \href{https://aimath.org/pastworkshops/geodesicsproblems.pdf}{[S1]}.

\subsection{Short English statement}
Do the lengths of closed geodesics, labelled by their loop classes, determine a negatively curved metric up to an isometry isotopic to the identity?
\subsection{Research attempt}
\textbf{Literal catalogue formulation disproved by an attributed high-dimensional
construction; standard homotopy-rigidity question unresolved here.}
\textbf{Scope and a formulation obstruction.}
The word ``isotopic'' in the supplied statement changes the problem in high
dimensions.  This section gives a counterexample to that literal conclusion,
conditional only on explicitly identified established topology theorems.  It
does not disprove the usual rigidity question with ``homotopic to the identity''
in the conclusion.  In the examples below an isometry homotopic to the identity
exists by construction.  The original statement above is retained so that this
difference remains visible.

The external input is the following consequence of Farrell--Jones [E1]: for a
closed connected real hyperbolic manifold $M$ of dimension $m>10$, the group
$\pi_0\operatorname{Diff}(M)$ of smooth isotopy classes is infinite.
The reference computes the infinite group $\pi_0\operatorname{Top}(M)$ in
Corollary 10.16 and discusses the stable splitting from diffeomorphisms in
equation (10.28).  Here compact support imposes no restriction because $M$ is
closed.  Butt [E2], Lemma 2.43 and Remark 2.44, explicitly records both the
finite quotient by diffeomorphisms homotopic to the identity and this infinite
quotient by diffeomorphisms isotopic to the identity.  We use only the real
hyperbolic case of the older theorem; no extension to all negatively curved
metrics is needed.  Closed real hyperbolic manifolds exist in these dimensions.
These topology and existence inputs are cited results, not proved by the
elementary argument that follows.

\textbf{An abstract pullback obstruction.}
Let $(M,g)$ be any closed Riemannian manifold for which the following hold:
\begin{enumerate}
\item the isometry group $I(g)$ is finite;
\item infinitely many smooth isotopy classes contain diffeomorphisms homotopic
to the identity.
\end{enumerate}
Choose a diffeomorphism $\phi$ homotopic to the identity whose isotopy class is
not the class of any member of $I(g)$.  There are only finitely many forbidden
classes, so the choice is possible.  Put $h=\phi^*g$.  Pullback preserves smoothness,
sectional curvature, and compactness; in particular, if $g$ has constant
curvature $-1$, so does $h$.

For any piecewise smooth loop $c$,
\[
 L_h(c)=L_g(\phi\circ c).
\]
Since $\phi$ is homotopic to the identity, the loops $c$ and $\phi\circ c$ have
the same free homotopy class.  Taking infima over that class, and using the
bijection of its loops induced by $\phi$, gives
\[
 \ell_h([c])=\ell_g([c]).
\]
Negative curvature ensures these infima are the lengths of the specified
unique geodesic representatives.  Thus the marking is preserved, not merely
the unmarked set of lengths.

Suppose the supplied conclusion held: there were $f$ smoothly isotopic to the
identity with $f^*h=g$.  Pullback reverses the order of composition, so
\[
 g=f^*(\phi^*g)=(\phi\circ f)^*g.
\]
Consequently $\phi\circ f\in I(g)$.  An isotopy from $f$ to the identity,
composed on the left with $\phi$, is an isotopy from $\phi\circ f$ to $\phi$.
This contradicts the choice of $\phi$.  Notice that selecting an arbitrary
nontrivial isotopy class in the homotopy kernel would not by itself suffice:
one must exclude the classes realized by isometries.  The finiteness/infinity
argument makes that exclusion explicit.

\textbf{Checking the two hypotheses for the hyperbolic example.}
Write
\[
 G=\operatorname{Diff}(M)/\operatorname{Diff}^{\rm iso}_0(M),\qquad
 H=\operatorname{Diff}^{\rm hom}_0(M)/\operatorname{Diff}^{\rm iso}_0(M).
\]
The first denominator consists of diffeomorphisms joined to the identity by a
smooth isotopy; the superscript ``hom'' means only homotopic as maps.
These are normal subgroups.  The quotient $G/H$ is finite by the cited
homotopy-class finiteness theorem.  If $H$ were finite, the finite union of its
cosets would make $G$ finite.  Farrell--Jones gives the opposite, so $H$ is
infinite.  This supplies the second hypothesis without identifying an explicit
coordinate formula for one of the exotic diffeomorphisms.

For completeness, finiteness of $I(g)$ follows from standard compact-isometry
group theory and a short curvature calculation.  The isometry group of a
closed Riemannian manifold is a compact Lie group.  An infinitesimal isometry
is a Killing field $X$.  The integrated Killing-field identity is
\[
 \int_M|\nabla X|^2\,dV=\int_M\operatorname{Ric}(X,X)\,dV.
\]
For a real hyperbolic $m$-manifold, the right side is
$-(m-1)\int_M|X|^2\,dV$.  Both sides can agree only when $X=0$.
The compact Lie group therefore has zero-dimensional Lie algebra, is discrete,
and hence is finite.  This argument uses the customary smooth isometry-group
theorem and the Bochner identity, neither a claim about arbitrary homeomorphisms
nor an assertion that every homotopy-trivial diffeomorphism is an isometry.

It follows that a closed real hyperbolic manifold in any dimension greater
than ten supplies a pair $(g,\phi^*g)$ satisfying the hypothesis and violating
the isotopy conclusion.  The inverse $\phi^{-1}$ is nevertheless homotopic to
the identity and obeys $(\phi^{-1})^*h=g$.  Therefore the very same example
satisfies the standard homotopy version of rigidity.

\textbf{What the remaining geometric rigidity question asks.}
The obstruction is not a failure to reconstruct lengths accurately.  Lengths
labelled by free homotopy classes are invariant under the entire group
$\operatorname{Diff}^{\rm hom}_0(M)$, so such data cannot distinguish its
different isotopy components.  One cannot repair the proof by extracting more
closed-orbit lengths from exactly the same marking.  A formulation using a
homotopy-trivial isometry, or an additional restriction making the two relevant
quotients agree, removes this particular obstruction.  Neither modification
is silently substituted for the inherited statement here.

There are useful elementary checks on potential arguments for the geometric
version.  If $h=a^2g$ with $a>0$, then every length is multiplied by $a$.  On a
closed negatively curved manifold there is a nontrivial closed geodesic, so
equal marked lengths force $a=1$.  Thus rescaling cannot furnish a competing
counterexample.  Equality of volume alone would be substantially weaker than
equality of every marked length and cannot replace it.

A second check concerns pointwise ordered metrics.  Suppose $h\geq g$ as
quadratic forms and their marked lengths agree.  If $c$ is the closed
$h$-geodesic in a nontrivial class, then
\[
 \ell_g([c])\leq L_g(c)\leq L_h(c)=\ell_h([c])=\ell_g([c]).
\]
Both inequalities are equalities.  The continuous nonnegative function
$|\dot c|_h-|\dot c|_g$ has integral zero, so it vanishes at every point of
this orbit.  Using the standard density of periodic geodesic-flow directions
for a closed negatively curved metric, it follows by continuity that
$(h-g)(v,v)=0$ on every tangent direction.  Polarization gives $h=g$.
Here periodic-direction density is an external dynamical theorem, and the
pointwise ordering is an extra hypothesis.  The pullback counterexamples need
not satisfy it; an arbitrary pair of metrics cannot be ordered by adding a
constant without changing its length spectrum.

\textbf{Linearization and why it does not remove the topology.}
For a differentiable family $g_t=g+tq+O(t^2)$ of negative-curvature metrics,
the first variation of the length in a fixed free homotopy class is
\[
 \left.\frac{d}{dt}\right|_{t=0}\ell_{g_t}([c])
 =\frac12\int_c q(\dot c,\dot c)\,ds,
\]
with $c$ parametrized at unit $g$-speed.  The curve-variation term vanishes
because $c$ is a closed geodesic.  If $q=\mathcal L_Vg$, then
$q(\dot c,\dot c)=2\frac{d}{ds}g(V,\dot c)$, whose integral is zero.
This proves the unavoidable infinitesimal gauge kernel directly.
It does not show that this is the entire kernel, and even an injectivity
theorem after fixing this gauge would address a neighborhood of a metric,
not all components of the diffeomorphism group.  The exotic pullbacks above
are a global isotopy-component obstruction, so an infinitesimal inverse
function theorem cannot by itself establish the supplied universal claim.

\textbf{Verification and outcome.}
The offline diagnostic checks pullback composition on positive definite
matrices and an elementary finite quotient example illustrating the coset
argument.  These are algebraic sanity checks only; the infinitude assertion
comes from the cited topology, not a finite computation.  The source audit
checked the original Farrell--Jones text and the explicit homotopy/isotopy
warning in Butt.  On those established inputs, the literal formulation is
disproved in high dimensions.  The general marked-length-spectrum rigidity
problem with a homotopy conclusion remains outside what this argument proves.

\subsection{Assessment of the claimed solution}

The argument above is mathematically sound as a counterexample to the
\emph{literal formulation stated in this document}, namely the formulation
requiring an isometry that is smoothly isotopic to the identity.  It should
not, however, be interpreted as a solution or disproof of the standard
Marked Length Spectrum Rigidity Conjecture.

The essential distinction is between being \emph{homotopic} to the identity
and being \emph{isotopic} to the identity.  Suppose that $(M,g)$ is a closed
real hyperbolic manifold of dimension greater than ten for which there exists
a diffeomorphism
\[
\phi:M\to M
\]
that is homotopic to the identity but whose smooth isotopy class is not
represented by any element of the finite isometry group $I(g)$.  The cited
high-dimensional topology results of Farrell--Jones [E1], together with the
homotopy/isotopy distinction recorded explicitly by Butt [E2], provide the
required supply of such isotopy classes.

Set
\[
h=\phi^*g.
\]
Because $\phi$ is homotopic to the identity, it preserves every free homotopy
class of loops.  Hence, for every nontrivial free homotopy class $[\gamma]$,
\[
\ell_h([\gamma])
=
\inf_{c\in[\gamma]} L_h(c)
=
\inf_{c\in[\gamma]} L_g(\phi\circ c)
=
\inf_{c\in[\gamma]} L_g(c)
=
\ell_g([\gamma]).
\]
Thus $g$ and $h$ have identical marked length spectra.

Now suppose that there were a diffeomorphism $f$ smoothly isotopic to the
identity such that
\[
f^*h=g.
\]
Then
\[
g
=
f^*(\phi^*g)
=
(\phi\circ f)^*g,
\]
so $\phi\circ f\in I(g)$.  Since $f$ is isotopic to the identity,
$\phi\circ f$ is isotopic to $\phi$.  This contradicts the choice of $\phi$,
whose isotopy class was assumed not to contain any $g$-isometry.  Therefore
the conclusion requiring an isometry isotopic to the identity fails.

However, this does \emph{not} produce two nonisometric metrics.  Indeed,
\[
(\phi^{-1})^*h
=
(\phi^{-1})^*(\phi^*g)
=
g,
\]
and $\phi^{-1}$ is homotopic to the identity.  Hence the two metrics are
isometric by an isometry preserving the marking up to homotopy.

This distinction is important because the usual marked length spectrum
rigidity problem is formulated in terms of a marking-preserving isometry,
equivalently an isometry homotopic to the identity, rather than necessarily
isotopic to the identity.  The AIM problem list [S1] asks whether equality of
the marked length spectra forces the metrics to be isometric, while [E2]
makes the relevant homotopy formulation explicit.

Consequently, the result established above should be described more
precisely as follows:

\begin{quote}
In dimensions greater than ten, the strengthened formulation of marked
length spectrum rigidity requiring the realizing isometry to be smoothly
isotopic to the identity is false.  The standard rigidity problem with a
homotopy-to-the-identity conclusion is not disproved by this construction.
\end{quote}

Accordingly, the label
\[
\texttt{ATTEMPT\_STATUS: solved}
\]
would be misleading if it is intended to refer to the standard Marked Length
Spectrum Rigidity Conjecture.  A more accurate description is:

\begin{quote}
\textbf{Counterexample to the isotopy-to-the-identity strengthening of
marked length spectrum rigidity in dimensions greater than ten; the standard
homotopy formulation remains unresolved by this argument.}
\end{quote}


\subsection{Sources}
\begin{itemize}
\item[S1]AIM workshop problem list, Open problems and questions about geodesics.
\url{https://aimath.org/pastworkshops/geodesicsproblems.pdf}.
\textit{Formulation source.}
\item[E1]F. T. Farrell and L. E. Jones, \emph{A Topological Analogue of Mostow's Rigidity Theorem}, Journal of the American Mathematical Society 2 (1989), 257--370; Corollary 10.16 and equation (10.28), pp. 350--351.
\url{https://math.uchicago.edu/~shmuel/tom-readings/FJ%20Hyperbolic.pdf}.
\textit{Published topology input; the argument uses closed real hyperbolic manifolds of dimension greater than ten.}
\item[E2]Karen Butt, \emph{Quantitative marked length spectrum rigidity}, Geometry \& Topology 29 (2025), Lemma 2.43 and Remark 2.44.
\url{https://arxiv.org/html/2203.12128}.
\textit{Explicit distinction between the homotopy quotient and the isotopy quotient of the diffeomorphism group.}
\end{itemize}
\end{document}
